\section{Introduction}
\label{sec:introduction}

Uniform spanning forests are canonical probability measures on the edges of
an infinite graph. Pemantle constructed infinite-volume spanning-tree limits
on the integer lattice~\cite[Theorem~2.3]{Pemantle1991}; Benjamini, Lyons,
Peres and Schramm developed their general graph theory~\cite{BLPS2001}. The free uniform spanning forest is obtained by taking
weak limits of uniform spanning trees on finite connected exhaustions, without
identifying boundary vertices. The limit exists and is independent of the
exhaustion. We recall its classical construction from finite tree determinants
in Proposition~\ref{prop:free-limit}; see also Benjamini, Lyons, Peres, and
Schramm~\cite[Section~5]{BLPS2001}.
A factor of IID description asks whether this global random object can be
constructed, equivariantly and measurably, from independent labels attached
to the vertices. Lyons explicitly discussed the FUSF factor question for
the group action on Cayley graphs in his report for the 2013 Oberwolfach
workshop~\cite[p.~2392]{Lyons2013}.

On transient graphs, the wired forest has an IID construction through
cycle popping and Wilson's algorithm rooted at infinity; see
\cite[Theorem~5.1]{BLPS2001}. Angel, Ray and Spinka give an explicit
graph-factor formulation, retaining the transient graph scope and crediting
the earlier construction~\cite[Section~4]{AngelRaySpinka2024}. Determinantal structure also provides
factor constructions in amenable group settings: Lyons and Thom prove
Bernoulli isomorphism for equivariant determinantal measures on amenable
Cayley graphs, with respect to the group action \cite[Corollary~7.4]{LyonsThom2016}. This Bernoulli-isomorphism result concerns
the specified group action; it does not by itself provide one rule
on arbitrary graphs equivariant under every graph isomorphism.

For general graph laws, Tim\'ar proves FIID representability of the free
forest on recurrent unimodular random graphs and, more generally, on
invariantly amenable unimodular random graphs
\cite[arXiv version, Theorem~1 and Corollary~5]{Timar2025}.
Under invariant amenability he obtains the stronger conclusion of a
finite-valued finitary factor~\cite[arXiv version, Theorem~2(1)]{Timar2025}.
His construction uses compatible monotone finite-volume limits along a
hyperfinite exhaustion. The introduction to that December 2025 version
states the unrestricted FUSF factor question as open. We prove the
unrestricted graphwise statement by a different construction: the same
Borel rule works for every graph in our class, independently of any
random-graph law. The challenge is to obtain a single equivariant sample
from the infinite-volume law while controlling all edges jointly.

\subsection{Statement}

The graphs in our main result are infinite, connected, locally finite, simple,
and undirected. Write $\FUSF_G$ for the free uniform spanning forest law on
$\{0,1\}^{E(G)}$. Graphs with a distinguished undirected edge and vertex
labels carry the usual Borel structure generated by finite marked
neighborhoods. A rule is \emph{root independent} if it uses no distinguished
vertex. Equivariance means invariance under every isomorphism of labeled
graphs carrying the distinguished edge to the distinguished edge.

\begin{theorem}\label{thm:main}
There is one Borel, root-independent, equivariant rule
\[
 \Phi(G,U,e)\in\{0,1\}
\]
such that, for every infinite connected locally finite simple undirected
unweighted graph $G$ and independent
$\operatorname{Uniform}[0,1]$ vertex labels $U=(U_v)_{v\in V(G)}$, the entire
edge set
\[
 \{e\in E(G):\Phi(G,U,e)=1\}
\]
has law $\FUSF_G$.
\end{theorem}

\begin{corollary}\label{cor:unimodular}
The free uniform spanning forest is a graph factor of IID under every
unimodular law supported on the stated class of graphs. The same rule works
for every such law.
\end{corollary}

Here unimodularity means the mass-transport identity
\[
 \E\sum_{v\in V(G)} f(G,o,v)
 =\E\sum_{v\in V(G)} f(G,v,o)
\]
for every nonnegative Borel function of a graph with two distinguished
vertices. Theorem~\ref{thm:main} is a graphwise statement and does not use
this identity. In particular, no amenability, transience, uniform degree
bound, or degree moment is assumed. The result asserts an ordinary Borel
factor, without a finitary coding radius or a finite input alphabet.

\subsection{Strongly Rayleigh processes}

The sampling criterion also applies to a broader class of binary laws.
For a finite set $F$, a probability measure $\mu$ on $\{0,1\}^F$ is
\emph{strongly Rayleigh} if its generating polynomial
\[
 g_\mu(z)=\sum_{x\in\{0,1\}^F}\mu(x)\prod_{i\in F}z_i^{x_i}
\]
is nonzero whenever $\operatorname{Im}z_i>0$ for every $i\in F$.
For a countable index set, we require every finite-coordinate marginal to
be strongly Rayleigh, following Borcea, Br\"and\'en, and
Liggett~\cite[Definition~2.13]{BorceaBrandenLiggett2009}.

For a countable group $\Gamma$, write $\tau_g$ for the regular left action
on coordinate families, given by $(\tau_g a)_h=a_{g^{-1}h}$.

\begin{theorem}[Strongly Rayleigh processes on countable groups]
\label{thm:sr-group}
Let $\Gamma$ be a countable group and let $\nu$ be a strongly Rayleigh
probability law on $\{0,1\}^{\Gamma}$ invariant under every $\tau_g$.
There is a Borel map
\[
 \Phi_\nu:[0,1]^{\Gamma}\longrightarrow\{0,1\}^{\Gamma}
\]
such that $\Phi_\nu(\tau_g u)=\tau_g\Phi_\nu(u)$ for every $g\in\Gamma$
and every $u\in[0,1]^{\Gamma}$, and such that $\Phi_\nu(U)$ has law $\nu$
when the coordinates of $U$ are independent $\operatorname{Uniform}[0,1]$
variables.
\end{theorem}

For a self-adjoint positive contraction $K$ on $\ell^2(\Gamma)$, the
determinantal law of a random subset $X\subseteq\Gamma$ is characterized
by the finite inclusion probabilities
\[
 \P(A\subseteq X)=\det(K_{g,h})_{g,h\in A}
 \qquad\text{for every finite }A\subseteq\Gamma.
\]

\begin{corollary}[Determinantal processes on countable groups]
\label{cor:dpp-group}
Let $\Gamma$ be a countable group and let $K$ be a self-adjoint positive
contraction on $\ell^2(\Gamma)$, so $0\le K\le I$. If the determinantal
probability law with kernel $K$ is invariant under the regular left action,
then it is a factor of IID under that action. In particular, this holds
whenever $K$ commutes with the left regular representation of $\Gamma$.
\end{corollary}

These statements concern the regular action on the group itself, not an
arbitrary action on another index set. The factor may depend on $\Gamma$
and $\nu$. No amenability or finite-generation assumption is imposed, and
finite groups are allowed. The conclusion is a Borel factor of IID; it
does not assert a finitary coding or a Bernoulli isomorphism. In particular,
it is distinct from the amenable-group Bernoulli-isomorphism result of
Lyons and Thom cited above. Their broader determinantal factor question
appears in~\cite[Question~7.7]{LyonsThom2016}; Corollary~\ref{cor:dpp-group}
treats the regular action on the group itself.

\subsection{The construction}

Let $X$ have the target forest law and observe its edge indicators through
independent Brownian noises:
\[
 Y_e(t)=tX_e+B_e(t).
\]
For finitely many observed edges, the posterior is an exponential tilt of
the forest law. The main quantitative input is the following response
identity for a finite weighted spanning tree. If $p_e$ is the probability
of including edge $e$ and $h_f$ is the logarithm of the weight of edge $f$,
then
\[
 \sum_f\left|\frac{\partial p_e}{\partial h_f}\right|
 =2p_e(1-p_e)\le\frac12.
\]
Negative off-diagonal covariances and the fixed size of a finite tree give
this identity. Passing it through finite-volume limits yields an
everywhere-defined posterior drift $b(t,y)$ whose response to bounded
changes of the field is uniformly Lipschitz, even when the fields themselves
are unbounded across the graph.

Subtracting the posterior drift gives the innovations
\[
 \widehat W_e(t)=Y_e(t)-\int_0^t b_e(s,Y(s))\,ds.
\]
They are independent Brownian motions. The response bound makes the integral
equation
\[
 Z_e(t)=w_e(t)+\int_0^t b_e(s,Z(s))\,ds
\]
uniquely solvable by a Borel Picard iteration for every continuous input
field $w$. Only differences between iterates are bounded; there is no
assumption that a countable Brownian field takes values in $\ell^\infty$.
Thus $Y=Z(\widehat W)$, and independent Brownian inputs to this deterministic
decoder reproduce the observation law. Integer-time slopes recover $X$.

Proposition~\ref{prop:brownian-sampler} isolates the resulting sampling
criterion for countable binary laws with uniformly controlled finite-tilt
responses. The criterion separates the probabilistic observation argument
from the quantitative property of the target measure. For spanning forests,
intrinsic edge-centered balls make the decoder equivariant across graphs.
Finally, independent vertex labels supply the independent edge paths by a
local allocation of disjoint random coordinates.

This construction adapts an established observation mechanism. Stochastic
localization was introduced by Eldan~\cite{Eldan2013}; a finite-dimensional
form with fixed noise covariance appears
in~\cite[Section~2, Proposition~7]{Eldan2020}. El Alaoui and
Montanari~\cite[Section~3, Theorem~2]{ElAlaouiMontanari2022} develop its
Gaussian-channel observation and innovation formulation. A particularly
close antecedent is Nam, Sly and
Zhang~\cite[Theorem~1 and Section~2.2]{NamSlyZhang2022}: for free Ising
measures on regular trees in their parameter range, they obtain FIID
codings through Brownian-driven systems with posterior-mean drift and
recover spins at large times. Here we prove the countable-coordinate
posterior and joint innovation identities, then use the uniform response
bound to obtain an inverse defined at every continuous input.

A companion~\cite[Theorem~1.1]{SharpFreeIsing2026} gives the sharp FIID
threshold for the free zero-field ferromagnetic Ising law on regular
trees, including the critical case, with a factor that can be chosen to commute
with every tree automorphism on every label input. Its conditional-copy recovery
differs from our uniform-response inverse.

For strongly Rayleigh laws, symmetric homogenization supplies a fixed-size
strongly Rayleigh extension, and the Rayleigh inequality gives nonpositive
off-diagonal covariances~\cite[Definition~2.12 and Theorems~4.1--4.2]{BorceaBrandenLiggett2009}.
The covariance-row bound then holds with the same constant $1/2$, as an
inequality rather than the weighted-tree identity. Translates of one finite
exhaustion of $\Gamma$ make the resulting drift equivariant for the regular
group action.

Section~\ref{sec:response} reconstructs the free limit and establishes the
response estimate and the canonical
drift. Section~\ref{sec:observation} proves the posterior and innovation
identities. The deterministic inverse and sampling criterion follow in
Section~\ref{sec:decoder}. Section~\ref{sec:implementation} constructs the
single Borel graph rule and proves Theorem~\ref{thm:main}.
Section~\ref{sec:strongly-rayleigh} proves Theorem~\ref{thm:sr-group} and
Corollary~\ref{cor:dpp-group}.
